%%% XeLaTeX-article %%%
%# -*- coding: utf-8 -*-
%!TEX encoding = UTF-8 Unicode
%!TEX TS-program = xelatex
%---------------------虽然加了%还是要保留！

\documentclass[17pt]{beamer}
\mode<presentation>
{
\usetheme[width=40pt]{Hannover}
\usecolortheme[]{dove}
\usefonttheme[]{structurebold}
\setbeameroption{hide notes}
}

\usepackage{fontspec}
\usepackage{polyglossia}
\setmainfont{Arial} %设置主字体
\newfontfamily\sanskritfont[Script=Devanagari,Mapping=romantodevanagari,Scale=1.15]{Sanskrit 2003}             %输出天城体
%\newfontfamily\sanskritfont[Mapping=tex-text]{Times New Roman}              %输出转写
\doublehyphendemerits=-10000
\newcommand{\skt}[1]{{\sanskritfont{#1}}} %输出天城体
\newcommand{\skttrans}[1]{{\skt{#1}~#1}}  %输出天城体和转写
%----------------------------------------------------设置梵文输入方法 danda । ॥

\usepackage[UTF8,fontset=windows]{ctex}
\usepackage{amsmath}
%----------------------------------------------------设置中文环境

\usepackage{graphicx}
\usepackage{flafter}
\graphicspath{{pic/}}
\usepackage{booktabs}
\usepackage{nicematrix}
\newenvironment{indentlist}
  {\begin{list}{}{\setlength{\leftmargin}{2em}\setlength{\itemsep}{0.5em}}}
  {\end{list}}
%-----------------------------------------插图表格

\usepackage{hyperref}
\usepackage[dvipsnames]{xcolor}
\usepackage{colortbl}
\definecolor{light-gray}{gray}{0.9}
%------------------------------颜色

\newcommand{\verbroot}[1]{\textcolor{red}{$\sqrt{}$#1}}
\newcommand{\sktroot}[1]{{\verbroot{\skt{#1}}}}
\newcommand{\skttransroot}[1]{{\sktroot{#1}~\textcolor{red}{#1}}}

\newcommand{\nounstem}[1]{\textcolor{red}{#1\nobreakdash-}}
\newcommand{\sktnounstem}[1]{{\textcolor{red}{\skt{#1\nobreakdash-}}}}
\newcommand{\skttransnounstem}[1]{{\sktnounstem{#1}~\nounstem{#1}}}

\newcommand{\verbstem}[1]{\textcolor{blue}{#1\nobreakdash-}}
\newcommand{\sktverbstem}[1]{{\textcolor{blue}{\skt{#1\nobreakdash-}}}}
\newcommand{\skttransverbstem}[1]{{\sktverbstem{#1}~\verbstem{#1}}}

\newcommand{\wordending}[1]{\textcolor{Orange}{\nobreakdash-#1}}
\newcommand{\sktending}[1]{{\textcolor{Orange}{\skt{-#1}}}}
\newcommand{\skttransending}[1]{{\sktending{#1}~\wordending{#1}}}

\newcommand{\fullpada}[1]{\textcolor{OliveGreen}{#1}}
\newcommand{\sktpada}[1]{{\textcolor{OliveGreen}{\skt{#1}}}}
\newcommand{\skttranspada}[1]{{\sktpada{#1}~\fullpada{#1}}}

\newcommand{\pratyaya}[1]{\mbox{\textcolor{Plum}{#1}}}
\newcommand{\sktpratyaya}[1]{{\textcolor{Plum}{\skt{#1}}}}
\newcommand{\skttranspratyaya}[1]{{\sktpratyaya{#1}~\pratyaya{#1}}}

\newcommand{\reconstruction}[1]{\textcolor{gray}{*#1}}
\newcommand{\fullsentence}[1]{\textcolor{MidnightBlue}{#1}}
\newcommand{\sktsentence}[1]{\textcolor{MidnightBlue}{\skt{#1}}}

\newcommand{\veryimportant}[1]{\textcolor{red}{#1}}
\newcommand{\important}[1]{\textcolor{blue}{#1}}
\newcommand{\notsoimportant}[1]{\textcolor{gray}{#1}}
\newcommand{\dicturl}[2]{\href{#1}{\mbox{\texttt{#2}}}}
%-------------------------------------------词根等标颜色

\title{{梵语提高}}
\subtitle{33. 愿望动词和必要分词}
\author[张雪杉]{文学院~~张雪杉 \\ zhangxueshan@sdnu.edu.cn}
\date{}

\begin{document}

\begin{frame}
  \titlepage
\end{frame}

\begin{frame}
  \frametitle{本节内容}
  \small
  \tableofcontents
\end{frame}

\section{上节作业}

\begin{frame}{第32章练习3}
  \small
  \begin{verse}
    \skt{1) tvayā vinā jīvituṃ necchāmītyuktvā kumāro 'pajagāma ।}\\
    \skt{2) tanmayā kṛtamiti hasantī bālāvadat ।}\\
    \mbox{\skt{3) sarvāṇi bhūtāni mayi vasantītīśvara uvāca ।}}\\
    \skt{4) ko yuvayoḥ śīghratara iti pṛṣṭau kumārau prativaktuṃ nāśaknutām ।}\\
  \end{verse}
\end{frame}

\begin{frame}{第32章练习3}
  \small
  \begin{verse}
    \skt{5) tanmama gṛhamidaṃ tu taveti naro mitrāyādarśayat ।}\\
    \skt{6) ratnāni no na santi । yuṣmākaṃ tu bahūni vasūnyeveti nārya ūcuḥ ।}\\
    \skt{7) ahaṃ tvanna kadā cidapagamiṣyāmi ।}\\
    \skt{8) kutaḥ sa vyāghro yuṣmābhirhata iti bālo 'pṛcchat ।}\\
  \end{verse}
\end{frame}

\begin{frame}{第32章练习3}
  \small
  \begin{verse}
    \skt{9) tvaṃ vardhethā āvābhyāṃ ca bahutarajño bhaveriti pitarau bālamavadetām ।}\\
    \skt{10) yuvābhyāṃ sahārīnabhibhavituṃ śakṣyāmaḥ ।}\\
    \skt{11) kuto 'smānna papracchitha । vayameva tvāṃ rakṣiṣyāma iti mitrāṇyūcuḥ ।}\\
    \skt{12) etanmamāsti tvayā tu luptam ।}\\
  \end{verse}
\end{frame}


\section{愿望动词}
\begin{frame}{\insertsection}
  \small
  \tableofcontents[currentsection]
\end{frame}

\subsection{意义}
\begin{frame}{\insertsection ~~意义}
  \begin{itemize}
    \item 愿望动词(desiderative)
    \vspace{1em}

    表示愿望、打算，\\或者行为应该发生。
  \end{itemize}
\end{frame}

\subsection{形态}
\begin{frame}{\insertsection ~~形态}
  \small
  \begin{itemize}
    \item 词根重复
    \item 加后缀 \pratyaya{-sa-}（有时要加联系元音 \pratyaya{-i-}，\\
      由于 \emph{ruki} 规则写作 \pratyaya{-ṣa-}、\pratyaya{-iṣa-}）
    \item 现在时语尾(语态与基础动词相同)
    \item \notsoimportant{愿望动词的完成时用迂回完成时（34课）}
  \end{itemize}
\end{frame}

\subsection{~~重复音节}
\begin{frame}{\insertsubsection}
  \small
  \begin{itemize}
    \item \textbf{辅音：}同完成时和第三类动词
    \item \textbf{元音：}
    \begin{itemize}
      \item 词根含 u/ū  → 重复元音 u
      \item 词根是其他元音 → 重复元音i
    \end{itemize}
  \end{itemize}
  \vspace{0.5em}
  \centering
  \resizebox{\textwidth}{!}{
    \begin{tabular}{@{}lll@{}}
      词根 & 愿望语干 & 含义 \\
      \midrule
      \verbroot{yudh} 'fight' & \verbstem{yuyutsa} (Ā.) & 'desires to fight' \\
      \verbroot{bhū} 'be' & \verbstem{bubhūṣa} (P.) & 'desires to be' \\
      \verbroot{pā} 'drink' & \verbstem{pipāsa} (P.) & 'is thirsty' \\
      \verbroot{bhṛ} 'carry' & \verbstem{bibhariṣa} (P.) & 'desires to carry' \\
    \end{tabular}
  }
\end{frame}

\subsection{~~不规则形式}
\begin{frame}{\insertsubsection}
  \small
  \begin{itemize}
    \item 词根元音 \pratyaya{-sa-} 前一般零级，\\
    \pratyaya{-iṣa-} 前一般二合，但总有例外：\\
    \verbroot{han} → \fullpada{jighāṃsati} 'wants to kill'
    \item 词根末 ṛ → \important{īr}：\verbroot{kṛ} → \fullpada{cikīrṣati}
    \item 短元音词根可延长：\verbroot{śru} → \fullpada{śuśrūṣate}
    \item 三个常见不规则形式：\\
    \verbroot{āp}  → \fullpada{īpsate} 'wants to get'\\
    \verbroot{dā}  → \fullpada{ditsati} 'wants to give'\\
    \verbroot{dhā} → \fullpada{dhitsati} 'wants to put'
  \end{itemize}
\end{frame}

\subsection{愿望形容词和名词}
\begin{frame}{愿望形容词}
  \small
  \begin{itemize}
    \item 将愿望动词语干末尾的 a 替换为 u\\
    \verbstem{yuyutsa} →  \nounstem{yuyutsu} 'desiring to fight' \\
    \verbstem{cikīrṣa} →  \nounstem{cikīrṣu} 'desiring to do' \\
    \verbstem{didṛkṣa} →  \nounstem{didṛkṣu} 'desiring to see' 
    \item 按以 u 结尾的形容词变格\\
    \item 例句 (MBh 2.11.2)\\
    \fullsentence{purā devayuge\ldots{} ādityo\ldots{} āgacchad}\\
    \fullsentence{mānuṣaṃ lokaṃ} \fullpada{didṛkṣuḥ}\\
    \notsoimportant{很久以前，在诸神的时代，太阳神}\\
    \notsoimportant{\underline{想要观看}人间，便来到此世。}
  \end{itemize}
\end{frame}

\begin{frame}{愿望名词}
  \small
  \begin{itemize}
    \item 将愿望语干末尾的 a 替换为 ā
    \item 按以 ā 结尾的阴性名词变格
    \item 一般表抽象含义
  \end{itemize}
  \begin{indentlist}
    \item \verbroot{pā} 'to drink' → \verbstem{pipāsa} \\
    → \nounstem{pipāsā} (f.) '渴（= 想要喝的欲望）'\\
    \item \verbroot{śru} 'to hear' → \verbstem{śuśrūṣa} \\
    → \nounstem{śuśrūṣā} (f.) '想要聆听；服从'\\
    \item \verbroot{bhuj} 'to enjoy' → \verbstem{bubhukṣa} \\
    → \nounstem{bubhukṣā} (f.) '饥饿'\\
  \end{indentlist}
\end{frame}


\section{必要分词}
\begin{frame}{\insertsection}
  \small
  \tableofcontents[currentsection]
\end{frame}

\subsection{意义}
\begin{frame}{\insertsection ~~意义}
  \small
  \begin{itemize}
    \item 必要分词(gerundive)\\
    也叫动形词或将来被动分词
    \item 由动词词根构成的\important{形容词}，\\
    表被动必要性 'having to be x-ed'
    \item 否定的必要分词（前缀 \pratyaya{a-}）\\译为 'must/should \emph{not} be \mbox{x-ed}' \\或 '\emph{cannot} be x-ed'
  \end{itemize}
\end{frame}

\subsection{形态}
\begin{frame}{\insertsection ~~形态}
  \small
  \begin{itemize}
    \item 词根（常二合，也有三合和零级）
    \item 后缀：\pratyaya{-(t)ya-} / \pratyaya{-tavya-} / \pratyaya{-anīya-}
    \item 格尾：阳性/中性按a变；阴性按ā变
  \end{itemize}
\end{frame}

\begin{frame}{三个后缀对比}
  \small
  \centering
  \begin{NiceTabular}{@{}lccc@{}}
    词根 & \pratyaya{-(t)ya-} & \pratyaya{-tavya-} & \pratyaya{-anīya-} \\
    \verbroot{kṛ} 'do' & \fullpada{kārya-} & \fullpada{kartavya-} & \fullpada{karaṇīya-} \\
    \verbroot{śru} 'hear' & \fullpada{śravya-} & \fullpada{śrotavya-} & \fullpada{śravaṇīya-} \\
  \end{NiceTabular}
  \begin{itemize}
    \item \pratyaya{-(t)ya-} 最常见，但实际都有可能。
    \item \pratyaya{-tya-} 是 \pratyaya{-ya-} 短元音后的变体
    \item \pratyaya{-tavya-} 前词根二合
    （基本与不定式 \pratyaya{-tum} 一致，有的要加联系元音\pratyaya{-i-} ）
    \item 不规则：词根末尾的 ā 在 \pratyaya{-ya-} 前变 e：\\
    \verbroot{dā} 'give' → \fullpada{deya-} 'having to be given'
  \end{itemize}
\end{frame}

\begin{frame}{例句}
  \small
  \sktsentence{sarve munayaḥ kuśīlavau} \sktpada{praśastavyau} \sktsentence{praśaṃsuḥ}\\
  \notsoimportant{所有圣者赞美了值得赞美的拘尸和罗婆。}\\
  \hfill\notsoimportant{(\emph{Rāmāyaṇa} 1.4.15)}
  \vspace{1em}

  \sktsentence{kairmayā saha} \sktpada{yoddhavyam} \sktsentence{asmin raṇasamudyame}\\
  \notsoimportant{在此战中，我必须与谁作战？}\\
  \hfill\notsoimportant{(\emph{Bhagavad-Gītā} 1.22)}
\end{frame}

\section{本节作业}

\begin{frame}{\insertsection }
  \begin{itemize}
    \item
      第33章练习1
  \end{itemize}
\end{frame}

\end{document}
